\documentclass[10pt,a4paper]{amsart}
\usepackage[margin=19mm]{geometry}
\usepackage{amsmath,amssymb,mathtools}
\usepackage[scr=rsfs,cal=euler]{mathalfa}
\usepackage{xfrac}
\usepackage[unicode,hidelinks,bookmarks=false]{hyperref}
\newcommand{\ie}{i.e.\ }
\newcommand{\cf}{cf.\ }
\newcommand{\resp}{respectively\ }
\newcommand{\Subsection}{\subsection}
\providecommand{\nobreakdash}{\nobreak\hbox{-}}
\setlength{\emergencystretch}{2em}
\multlinegap0pt
\usepackage{tikz}
\usepackage{mathtools}
\usepackage{amssymb}
\usepackage{xcolor}
\usepackage{algorithm}\usepackage{algpseudocode}
\newtheorem{lemma}{Lemma}
\newtheorem{remark}{Remark}
\newcommand{\boldc}{\mathbf{c}}
\newcommand{\boldf}{\mathbf{f}}
\newcommand{\boldq}{\mathbf{q}}
\newcommand{\bolds}{\mathbf{s}}
\newcommand{\boldu}{\mathbf{u}}
\newcommand{\boldv}{\mathbf{v}}
\newcommand{\boldz}{\mathbf{z}}
\newcommand{\cC}{\mathcal{C}}
\newcommand{\cS}{\mathcal{S}}
\title{Decoding From One Fragment in Forensic Printing\\Using Van der Corput sets}
\author{Junsheng Liu, Netanel Raviv}
\date{Source: arXiv:2609.00522v2 (CC BY 4.0)}
\begin{document}
\maketitle

\noindent\emph{Source and licence.} Decoding From One Fragment in Forensic Printing\\Using Van der Corput sets. Version arXiv:2609.00522v2 (CC BY 4.0), distributed by arXiv under the Creative Commons Attribution 4.0 International licence, http://creativecommons.org/licenses/by/4.0/, as recorded in the arXiv licence metadata for this exact version and shown on the abstract page https://arxiv.org/abs/2609.00522v2. Full licence notice: \texttt{licenses/arxiv-2609.00522-CC-BY-4.0.txt}.\par\smallskip

\noindent\textit{Editorial scope.} Counted unit: the article's lemma lemma:recycled-fragment-ambiguity (source0.tex lines 1833--1844). The article's complete original proof of it is delivered with it (source0.tex lines 1846--2065, one \texttt{proof} environment of 220 lines). No mathematics is written, repaired or re-derived: the statement and the proof are the article's own lines, in the article's own notation, and no retained line is substituted or re-flowed.  Retained uncounted as the source-local context the proof needs: lines 1720--1725; lines 1726--1731; lines 3400--3400; lines 3515--3523.  Nothing was omitted from any retained span, so the package carries no omission record.  No retained line contains a citation, so the wrapper carries no bibliography and no bibliography entry.  No retained line contains a figure environment, an image-inclusion command or drawing code, so no image is delivered and the package declares no asset dependency.  Editorial formatting only, in addition: an AMS \texttt{amsart} preamble that restates the article's own declarations and macros used by the retained lines, namely the environments \texttt{lemma}, \texttt{remark} on separate counters, together with the standard packages those lines need.  The retained environments print in this document as Lemma 1, Remark 1 in order of appearance, whereas the article prints the same items with the numbering of its own full document.  The complete native source of the article as distributed in the arXiv e-print for this exact version is bundled unchanged as \texttt{sources/209036/source0.tex}.
\begin{lemma}\label{lemma:usable-level-signature}
Fix $i\in[J-1]$, assume that levels $i+1,\ldots,J$ have been decoded correctly, i.e., every bit in $\cS_i$ is known to the decoder.
Then, for all sufficiently large $n$, every fragment of length at least
$(1+\delta)t_i\log n$ contains at least $(1+\delta/2)(t_i/T)\log n$ bits from the fixed De Bruijn sequence and at least $(1+\delta/2)(1-t_i/T)\log n$ bits from the already decoded level codewords $\bar{\bolds}^{\,i+1},\ldots,\bar{\bolds}^{\,J}$, where
$\bar{\bolds}^{\,j}\in\bar{\cC}_{\mathrm{er}}^j$ for all~$j\in\{i+1,\ldots,J\}$.
\end{lemma}

\begin{remark}\label{remark:special bits in S_j}
For the latter count $(1+\delta/2)(1-t_i/T)\log n$ in Lemma~\ref{lemma:usable-level-signature}, we only count the $\beta-1$ positions in each
length-$\beta$ RLL block that come from the shifted codeword in
$\cC_{\mathrm{sh}}^j$, $j\in\{i+1,\ldots,J\}$, and do not count the fixed $1$ inserted by the $\mathrm{RLL}(0,\beta-1)$ encoder.
Furthermore, it follows from Lemma~\ref{lemma:usable-level-signature} that every fragment of length at least $(1+\delta)t_i\log n$ contains at least $(1+\delta/2)\log n$ known bits of the two types"
\end{remark}

\section{Supporting estimates for rate and reliability}\label{appendix:rate-reliability-estimates}

\begin{lemma}\label{lemma:number-of-debruijn-subsequences-in-fragment}
Fix $i\in[J-1]$ and consider a fragment $\boldf$ of length
$(1+\delta)t_i\log n$.
Consider all candidate length-$k_p$ De Bruijn subsequences that can be
extracted from $\boldf$ under all possible candidate placements.
The number of such candidate subsequences, counted with multiplicity
according to how they arise from the fragment, is at most
\(Tb_n(1+\delta)t_i\log n=n^{o(1)}\).
\end{lemma}

\begin{lemma}\label{lemma:recycled-fragment-ambiguity}
For every fixed \(i\in[J-1]\), there exists a constant \(\rho_i>0\) such that for every fragment $\boldf=(\boldf[0],\boldf[1],\ldots,\boldf[(1+\delta)t_i\log n-1])$ of length \((1+\delta)t_i\log n\), the probability that there exist two or more $i$-consistent placements 
is at most $n^{-\rho_i+o(1)}$.
In particular, the constant \(\rho_i\) may be chosen to satisfy
\[
0<\rho_i<
\min\left\{
\frac{\delta}{2},
\left(1+\frac{\delta}{2}\right)\left(1-\frac{t_i}{T}\right)
\right\}.
\]
\end{lemma}

\begin{proof}
Let the true starting position for a fragment~$\boldf$ of length $(1+\delta)t_i\log n$ be \(s\in\{0,1,\ldots,n-1\}.\)
By Lemma~\ref{lemma:usable-level-signature}, every candidate placement of
length \((1+\delta)t_i\log n\) contains at least
\(k_{p}=(1+\delta/2)(t_i/T)\log n\) fixed De Bruijn bits and
at least
\(k_{r}=(1+\delta/2)(1-t_i/T)\log n\) bits from $\bar{\bolds}^{\,i+1},\ldots,\bar{\bolds}^{\,J}$; see Remark~\ref{remark:special bits in S_j}.

Suppose that there exists at least one false placement \(a \neq s\) that is also consistent.
It suffices to bound the probability of the union of these events over all \(a \neq s\).
For a false placement to be consistent, every fragment bit mapped to one of the known positions (i.e.,~$\cS_i$ minus the RLL bits and marker bits, see Remark~\ref{remark:special bits in S_j}) must agree with the known bit at that position. 
We first observe the randomness of the level-codeword bits used in these
comparisons. Fix the message and all parity-check matrices. For every
level \(h\in[J]\), we have
\(\boldv^{\,h}=\boldu^{\,h}\oplus\boldz^{\,h}\), where
\(\boldz^{\,h}\) is uniform over
\(\mathbb F_2^{n_{\mathrm{er}}^h}\), and the shifts are independent
across different levels. Therefore, the coordinates of
\(\boldv^{\,1},\ldots,\boldv^{\,J}\) are mutually independent
\(\operatorname{Bernoulli}(1/2)\) random variables. Moreover, by
Remark~\ref{remark:special bits in S_j}, each of the~\(k_r\) positions
counted in Lemma~\ref{lemma:usable-level-signature} contains a bit from one of the shifted level codewords, rather than a marker bit or
a~\(1\) added by the RLL encoder. Such distinct codeword positions
correspond to distinct coordinates of the shifted level codewords.

Our proof strategy is as follows. 
For a false placement disjoint from the true placement, matching with $k_{p}$ fixed De Bruijn bits reduces the number of possible false positions, while matching with the $k_{r}$ bits from the already decoded level codewords reduces the probability of that placement.
%proves such false placement happens with low probability. 
For a false placement overlapping the true placement---of which there are at most \(3(1+\delta)t_i\log n\)---$k_{r}$ bits from the already decoded level codewords alone are used to prove that such false placement is unlikely to happen.

We first consider false placements disjoint from the true placement, namely
\( |a-s|\geq(1+\delta)t_i\log n \).
By Lemma~\ref{lemma:usable-level-signature}, the fragment contains at least
\(k_p\) bits from the fixed De Bruijn sequence.
The decoder does not need to know which bits of the fragment are from the
fixed De Bruijn sequence.
Rather, once a candidate \(a\) is chosen, the alignment determines which
fragment indices
\(\boldf[j]\), \(j\in\{0,1,\ldots,(1+\delta)t_i\log n-1\}\),
would be placed at the fixed De Bruijn sequence positions.
For the candidate to be consistent, all of these observed bits
\(\boldf[j]\) must agree with the corresponding known fixed De Bruijn bits.

Furthermore, since the De Bruijn sequence is of order
\(\log N_{q,n}\), any fixed binary sequence of length \(k_p\) occurs
in \(\boldq\) at most
\(2^{\max\{\log N_{q,n}-k_p,0\}}\) times.
Indeed, if \(k_p\leq\log N_{q,n}\), then the length-\(k_p\) sequence
has at most \(2^{\log N_{q,n}-k_p}\) possible extensions to a
length-\(\log N_{q,n}\) sequence, and each such extension occurs once
in the De Bruijn sequence. If \(k_p>\log N_{q,n}\), then the sequence
can occur at most once, since its first \(\log N_{q,n}\) bits already
determine its starting position uniquely. These two cases are expressed
together by the exponent
\(\max\{\log N_{q,n}-k_p,0\}\).

Now fix one of these remaining disjoint false placements and select
\(k_r\) of its comparisons against known bits from the already decoded
level codewords. Since the false placement is disjoint from the true
placement, the codeword positions used as the known bits in these
\(k_r\) comparisons lie outside the positions occupied by the fragment
under its true placement. Hence, by the one-to-one correspondence
between these codeword positions and coordinates of the shifted level
codewords, none of the corresponding random-shift coordinates occurs
among the observed fragment bits.


Condition on the message, all parity-check matrices, the channel cuts,
the entire fragment \(\boldf\), and all random-shift coordinates except
the \(k_r\) distinct coordinates corresponding to these known
level-codeword bits. Under this conditioning, whether the fixed
candidate passes the De Bruijn comparisons is already determined.
The \(k_r\) remaining level-codeword bits are still mutually independent
\(\operatorname{Bernoulli}(1/2)\) random variables, while the
corresponding fragment bits are fixed. Therefore,
\[
\Pr\left(
\text{all \(k_r\) fragment bits aligned against level-codeword bits are consistent}
\right)
\leq 2^{-k_r}.
\]
Since this bound holds after conditioning on the fragment and on the
outcome of the De Bruijn comparisons, it also holds without this
conditioning.

According to
Lemma~\ref{lemma:number-of-debruijn-subsequences-in-fragment} in Appendix~\ref{appendix:rate-reliability-estimates}, there are
at most \(n^{o(1)}\) candidate length-\(k_p\) De Bruijn subsequences in
a segment of length \((1+\delta)t_i\log n\).
Applying a union bound over all disjoint false placements that survive
the De Bruijn comparisons gives
\[
\Pr\left(
\text{there exists a consistent disjoint false placement}
\right)
\leq
n^{o(1)}
2^{\max\{\log N_{q,n}-k_p,0\}-k_r}.
\]
We now consider the two cases in the maximum.
On the one hand, if \(k_p\leq\log N_{q,n}\), then
\[
\max\{\log N_{q,n}-k_p,0\}-k_r
=
\log N_{q,n}-k_p-k_r.
\]
Since
\(k_p+k_r=(1+\delta/2)\log n\) by Lemma~\ref{lemma:usable-level-signature} and since \(\log N_{q,n}=\log n+O(1)\), it follows that
\[
\log N_{q,n}-k_p-k_r
\leq
-\frac{\delta}{2}\log n+O(1).
\]
Hence, in this case,
\[
\Pr\left(
\text{there exists a consistent disjoint false placement}
\right)
\leq
n^{-\frac{\delta}{2}+o(1)}.
\]
On the other hand, if \(k_p>\log N_{q,n}\), then the maximum equals \(0\), and therefore
\[
\Pr\left(
\text{there exists a consistent disjoint false placement}
\right)
\leq
n^{o(1)}2^{-k_r}
=
n^{-\left(1+\frac{\delta}{2}\right)
\left(1-\frac{t_i}{T}\right)+o(1)}.
\]
Thus, in both cases,
\begin{equation}\label{equation:disjoint false placement}
    \Pr\left(
\text{there exists a consistent disjoint false placement}
\right)
\leq
n^{-\min\left\{
\frac{\delta}{2},
\left(1+\frac{\delta}{2}\right)
\left(1-\frac{t_i}{T}\right)
\right\}+o(1)}.
\end{equation}



We next consider a fixed false placement \(a\neq s\) that overlaps the
true placement, and let \(d\triangleq a-s\). Each of the~\(k_r\)
comparisons against previously decoded level-codeword bits induces an
equality constraint
\[
\boldc[s+j]=\boldc[s+j+d]
\]
between two codeword coordinates whose indices differ by \(d\). By
Remark~\ref{remark:special bits in S_j}, the coordinate
\(\boldc[s+j+d]\) in each such comparison is an original bit from one
of the shifted level codewords, rather than a marker bit or a \(1\)
added by the RLL encoder.
Construct a graph whose vertices are the codeword coordinates involved
in these comparisons and whose edges represent these equality
constraints, with each edge directed from \(s+j\) to \(s+j+d\).
Since every edge has displacement \(d\), each connected component lies
within a single residue class modulo \(|d|\) and is a subgraph of a
path. Hence, the comparison graph is a forest, or more specifically, a
collection of paths.

Consider one connected component containing \(k\) edges. With the above
orientation, every vertex in the component except for one endpoint is
the head of exactly one edge. Therefore, these \(k\) vertices are all
positions containing original shifted level-codeword bits. Moreover,
they correspond to \(k\) distinct coordinates of the random shifts and
are therefore mutually independent
\(\operatorname{Bernoulli}(1/2)\) random variables. The remaining
endpoint may be either a shifted level-codeword bit or a deterministic
pilot or an RLL bit, which does not affect the argument.
Conditioning on the value of this remaining endpoint, the \(k\)
equality constraints along the path successively require the \(k\)
independent \(\operatorname{Bernoulli}(1/2)\) bits to take prescribed
values. Hence, all \(k\) constraints in this component are satisfied
with probability at most \(2^{-k}\). Since different connected
components involve distinct shifted-codeword coordinates, multiplying
over all components gives
\[
\Pr\bigl(
\text{a fixed overlapping false placement is consistent}
\bigr)
\leq
2^{-k_r}.
\]

A union bound over the \(O(\log n)\) possible overlapping false placements gives
\begin{equation}\label{equation:overlap false placement}
    \Pr\left(
\text{there exists a consistent overlapping false placement}
\right)\leq
3(1+\delta)t_i\log n\,2^{-k_{r}}=
n^{-
\left(1+\frac{\delta}{2}\right)
\left(1-\frac{t_i}{T}\right)+o(1)}.
\end{equation}


%Combining the disjoint and overlapping false placements, we obtain
By adding~\eqref{equation:disjoint false placement} and~\eqref{equation:overlap false placement} we obtain the following via union bound:
\[
\Pr\bigl(\text{there exists at least two $i$-consistent placements})
\leq
n^{-\rho_i+o(1)}
\]
for every constant
\[
0<\rho_i<
\min\left\{
\frac{\delta}{2},
\left(1+\frac{\delta}{2}\right)
\left(1-\frac{t_i}{T}\right)
\right\}.\qedhere
\]
\end{proof}

\end{document}
